\section{Introduction}\label{sec:intro}

Process models tend to be wrong precisely where constraints bind. This
paper works through a case in which the wrongness can be put on a
scale. The classical kinetic law for amorphous-silica polymerization,
transcribed faithfully from its 1979 primary source, mispredicts the
effective decay rate observed in an independent 2024 laboratory corpus
by a factor of six at the median. Yet silica scaling must be managed
daily and at industrial scale. In the liquid-dominated geothermal
fields of the Philippines, which rank among the largest installed
capacities in the world~\citep{fronda2020country}, dissolved silica
sets the floor on reinjection temperature and therefore the ceiling on
bottoming-cycle power. Mitigation practice at the Bacman- and
Tiwi-type fields that motivate this study is documented and static:
fixed inhibitor programs and manually managed reinjection
temperatures~\citep{see2010wgc,remoroza2008ph,baltazar2014grc,%
kobayashi2024mahanagdong}. Between a model too wrong to trust and a
practice too static to optimize sits the control problem addressed
here.

The individual ingredients are established. Machine-learned predictors
of geothermal scaling exist in prediction-only
form~\citep{juhri2024sgw}. Hybrid physics--data modeling has a mature
canon~\citep{psichogios1992hybrid,thompson1994modeling,%
vonstosch2014hybrid,karniadakis2021physics}, soft-sensor practice under
sparse labels is well codified~\citep{kadlec2009data}, and real-time
optimization under plant--model mismatch has a developed theory of
modifier adaptation~\citep{marchetti2009modifier,marchetti2010dual,%
darby2011rto}. What has been missing, and what a constraint with
fourth-power sensitivity punishes, is the coupling between them: an
estimator that knows how much of its own learned correction to believe,
and an optimizer that spends that belief as constraint margin. Two
pre-registered negative results in this paper show that the coupling is
not a refinement but a requirement. A residual transferred across
plants actively harms, and a scalar bias adaptation, the natural first
recourse, cannot stabilize a multiplicative kinetic error on a
fourth-power constraint even while it harvests spectacular nominal
profit.

The contributions are fivefold. First, a credibility-weighted hybrid
soft sensor whose leave-one-out shrinkage yields provable graceful
degradation to its physics prior and doubles as the deployment safety
interlock. Second, a vertex-robust economic real-time optimizer that
consumes the sensor's correction together with its cross-validated
uncertainty as a multiplicative back-off. Third, a two-regime
closed-loop evaluation, spanning benign random mismatch and an adverse
stress set at the measured sixfold gap, that maps a safety--harvest
frontier on which the proposed loop is the only policy essentially free
of violations in both regimes. Fourth, out-of-sample validation of
every physics component and of the sensing pattern itself against the
independent 2024 corpus, which yields a reconciliation of second-order
decay with the classical fourth-order peak-rate law, a mutual
solubility validation at a median of $1.05\times$, and a design
principle we state as: identifiability precedes correction. Fifth, a
fully specified and regenerable benchmark problem, built on
Philippine-field chemistry, with a quartic constraint, an induction
delay, a non-monotonic actuator and a $\times 101$ temperature lever,
offered to the constrained-optimization community.

Section~\ref{sec:benchmark} states the problem and its industrial
context, Section~\ref{sec:arch} the architecture,
Section~\ref{sec:results} the in-silico results,
Section~\ref{sec:validation} the out-of-sample validation, and
Section~\ref{sec:discussion} the limitations and what generalizes.
